\documentclass[12pt]{article}
\usepackage{siunitx}

\begin{document}

\section*{Internal Team Meeting Agenda: Metrics \& Project Alignment}

\subsection*{Meeting Objectives}
\begin{itemize}
    \item Align the team on \textbf{what exactly we are testing} in the project
    \item Define \textbf{clear, measurable metrics} (KPIs) for both software and hardware components
    \item Establish \textbf{minimum passing criteria and failure conditions}
    \item Identify \textbf{existing benchmarks and prior work} to ground our evaluation
    \item Agree on a concrete \textbf{short-term action plan} for the next week
\end{itemize}

\subsection*{1. Project Goal Re-alignment}
\begin{itemize}
    \item Restate the core goal of the second semester project
    \item Clarify a measurable goal for the overall system:
    \begin{itemize}
      \item Objective from first semester:
      To demonstrate successful end-to-end handover, the system must integrate language understanding, symbolic planning, perception, and grasping into a reliable pipeline. Recent studies show that robot-to-human tool handover can reach around 92.5\% success in simulation for construction tools~\cite{iaarc2025_handover}, while sim-to-real grasping frameworks report 90--97\% success depending on object familiarity~\cite{mogpe2022},\cite{grasping2023}. Performance is typically lower for novel or cluttered settings, highlighting the importance of generalization.  

      Timing benchmarks suggest that simple object handovers can be completed in 8--10 seconds~\cite{handover2024_fast}, though more complex engineering tools may reasonably require up to 15 seconds. Based on these findings, this project defines its main performance goal as achieving \textbf{$\geq90\%$ success over 60–80 trials}, with \textbf{completion within 15 seconds} and \textbf{minimal disturbance} $\leq\SI{2}{\centi\meter}$ to non-target tools.
    \end{itemize}

\end{itemize}

\subsection*{2. Define the Final System-Level Test}
\begin{itemize}
    \item Describe the \textbf{final end-to-end task} the system must perform
    \item Identify real-world inspired task scenarios to benchmark against
    \item Discuss:
    \begin{itemize}
        \item What success looks like
        \item What constitutes a clear failure
    \end{itemize}
    \item Decide on observable outputs for evaluation
\end{itemize}

\subsection*{3. Metrics and KPIs Definition}
\textbf{Discussion focus: be specific, measurable, and benchmarked}
\begin{itemize}
    \item Performance metrics:
    \begin{itemize}
        \item Task completion time
        \item Planning latency
        \item Execution success rate
    \end{itemize}
    \item Accuracy metrics:
    \begin{itemize}
        \item Vision perception accuracy
        \item Correctness of generated PDDL
    \end{itemize}
    \item Robustness metrics:
    \begin{itemize}
        \item Failure recovery rate
        \item Sensitivity to perception errors
    \end{itemize}
    \item Define:
    \begin{itemize}
        \item Minimum passing criteria
        \item Failing conditions
    \end{itemize}
\end{itemize}

\subsection*{4. Component-Level Evaluation Plan}
\begin{itemize}
    \item Vision module evaluation
    \item Ontology-based PDDL generation evaluation
    \item Planner integration with vision outputs
    \item Decide what to test independently vs. end-to-end
\end{itemize}

\subsection*{5. Benchmarks and Prior Work}
\begin{itemize}
    \item Identify existing benchmarks or datasets to reuse
    \item Review prior work on:
    \begin{itemize}
        \item Automatic PDDL generation
        \item Ontology-driven task planning
        \item Vision-planning integration
    \end{itemize}
    \item Decide which benchmarks are realistic and defensible
\end{itemize}

\subsection*{6. One-Week Brainstorming Phase (No Development)}
\begin{itemize}
    \item Agree on a one-week period focused on:
    \begin{itemize}
        \item Metric refinement
        \item Literature review
        \item Benchmark selection
    \end{itemize}
    \item Define expected outputs at the end of the week
\end{itemize}

\subsection*{7. Action Items and Ownership}
\begin{itemize}
    \item Assign responsibilities for:
    \begin{itemize}
        \item Metric definition
        \item Benchmark research
        \item Ontology and PDDL review
        \item Vision-planning integration analysis
    \end{itemize}
    \item Set internal deadlines and next meeting date
\end{itemize}

\subsection*{8. Wrap-up}
\begin{itemize}
    \item Confirm shared understanding of goals and metrics
    \item Identify open questions or risks
    \item Align on next steps
\end{itemize}

\newpage
\bibliographystyle{IEEEtran}
\bibliography{ref}

\end{document}